\section{How data centers run LLM clouds}\label{sec:datacenter}

Hyperscalers and GPU-cloud operators serve models on large, uniform fleets run by a cluster
scheduler (Google's Borg and its open descendant Kubernetes), with inference layers such as Ray,
the vLLM production stack, llm-d and NVIDIA Dynamo on top. Ymir solves a much smaller version of
the same problem, a handful of mismatched machines owned by one person, so it borrows the ideas
that still pay off at that size and leaves out the ones that only make sense for a fleet.

{\small
\begin{xltabular}{\linewidth}{@{}L{2.7cm}YY@{}}
\toprule
Concern & Typical data center & Ymir \\
\midrule
\endhead
Scheduling & Cluster scheduler (Borg, Kubernetes) places containers; Ray or an operator manages model replicas & No cluster scheduler; one inference node will hold a lease and write the placement plan (designed, M6) \\
Load balancing & KV-cache and load-aware routers (llm-d, Dynamo, Gateway API Inference Extension) & Power-of-two-choices on predicted completion, from signed capability records \\
Prefill and decode & Often disaggregated onto separate GPU pools (DistServe, Dynamo, llm-d) & Not disaggregated; each engine does both \\
Autoscaling and placement & Replicas scale with traffic and SLOs; new capacity is provisioned & Fixed hardware; designed replication and cross-node swap within what you own \\
Failure handling & Retry on another replica; Dynamo can migrate a request with its tokens & Resume from the receiver's buffer with offsets, at model and coordinator level \\
Tensor parallel networking & InfiniBand or NVLink fabrics, RDMA & Opt-in split groups on a declared private link (cable or isolated VLAN), never over Weirdnet \\
Tenancy & Many tenants, quotas, rate limits, billing & One owner; fairness only between devices and priority classes \\
Edge auth & Public HTTPS gateway with API keys or OAuth & No public endpoint; Noise-authenticated peers plus the identity roster \\
Observability & Prometheus, tracing, SLO dashboards across the fleet & \code{ymirctl status}, optional Prometheus text on loopback, alerts by email and webhook \\
Hardware & Homogeneous accelerators per pool & Heterogeneous personal hardware; capabilities advertised per node \\
\bottomrule
\end{xltabular}
}

\noindent Ideas Ymir borrows, each from its design doc:
\begin{enumerate}
  \item Power of two choices: score two random candidates and take the better, so stale load data does not cause herding (Mitzenmacher).
  \item Priority classes: interactive, batch and background, with batch yielding between requests.
  \item Admission control at the server: model nodes answer Busy when slots and queue are full, so several coordinators never overcommit a node.
  \item Leases for a single writer: placement is decided by one inference node holding a renewable lease, while every node keeps routing (designed).
  \item Signed, expiring capability records, much like heartbeat and node-status reporting in cluster managers.
  \item Prefix-aware routing as a cost reduction, not a requirement (designed, from synced prefix hints).
  \item Draining before maintenance, and per-device fair sharing in the spirit of deficit round robin.
\end{enumerate}

\noindent What Ymir deliberately does not do:
\begin{enumerate}
  \item Run a general cluster scheduler, containers or Kubernetes.
  \item Disaggregate prefill and decode, or move KV caches between nodes; KV stays local and is rebuilt or warmed.
  \item Autoscale by buying or provisioning capacity; the pool is the hardware you own.
  \item Carry tensor-parallel traffic over its network; that stays on a fast private link.
  \item Serve other tenants, issue API keys, enforce quotas or expose a public gateway.
\end{enumerate}

The result is a design that looks like a miniature inference cloud from the inside (routing,
admission, priorities, failover, synced state) but is shaped for trust and heterogeneity rather
than for scale.

\subsection{At a glance: Ymir against the usual home stack}\label{sec:comparison-summary}

{\small
\begin{xltabular}{\linewidth}{@{}L{2.4cm}YYYY@{}}
\toprule
 & Ymir & llama-swap & Hosted APIs & Single-node servers \\
\midrule
\endhead
Scope & Pool of your own machines & One machine & Provider's cloud & One machine \\
Remote access & Weirdnet only, roster-checked & HTTP; add proxy or VPN & Public HTTPS & HTTP port; add proxy or VPN \\
Credentials & Device identity; no API keys & Optional keys & API keys & Usually none or proxy keys \\
Model swapping & Per node now; pool-wide designed & Per machine & Provider-managed & Per machine \\
Mid-answer failover & Yes, both levels & No & Provider-dependent & No \\
Where data goes & Only your linked devices & Your machine & Provider & Your machine \\
\bottomrule
\end{xltabular}
}

llama-swap is a per-machine proxy that starts and stops model servers on demand; Ymir absorbs that
job into each model node's supervisor and adds routing across machines, so it replaces llama-swap
rather than wrapping it. Hosted APIs offer scale and frontier models Ymir does not aim to match,
at the cost of sending data off your hardware and managing keys. Single-node servers in the
Ollama style are the simplest way to run a model on one box; Ymir's one-box mode covers the same
ground, and stays OpenAI-compatible on loopback for local tools, but its point is what happens
when you add the second machine.
